\section{Limitations}

ROK is a structural reliability mechanism, not a correctness guarantee.

\subsection{No Guarantee of Correctness}
Adversarial critique improves failure visibility but cannot eliminate undetected errors, especially in specialized domains.

\subsection{Dependence on Critic Quality and Policy Design}
Effectiveness depends on critic strength and decision policy thresholds. Weak critique or permissive thresholds reduce interception.

\subsection{Latency and Compute Cost}
Mandatory critique and bounded revision add latency and cost relative to single-agent pipelines.

\subsection{Open-Core Execution Constraints}
The reference implementation excludes side-effectful execution, limiting direct applicability without careful integration.

\subsection{Scope of Failure Coverage}
ROK targets failures rooted in assumptions, objective collapse, and premature convergence, but is less effective against incomplete world models or adversarial inputs without external validation.

\subsection{Not a Substitute for Human Judgment}
ROK does not replace human oversight; it improves inspectability and surfaces issues earlier in the pipeline.

\subsection{Summary}
ROK’s value lies in transforming silent failure into explicit, auditable structure.

